\documentclass[10pt,a4paper]{amsart}
\usepackage[margin=19mm]{geometry}
\usepackage{amsmath,amssymb,mathtools}
\usepackage[scr=rsfs,cal=euler]{mathalfa}
\usepackage{xfrac}
\usepackage[unicode,hidelinks,bookmarks=false]{hyperref}
\newcommand{\ie}{i.e.\ }
\newcommand{\cf}{cf.\ }
\newcommand{\resp}{respectively\ }
\newcommand{\Subsection}{\subsection}
\providecommand{\nobreakdash}{\nobreak\hbox{-}}
\setlength{\emergencystretch}{2em}
\multlinegap0pt
\usepackage{bm}
\usepackage{amssymb}
\usepackage{mathtools}
\usepackage[shortlabels]{enumitem}
\usepackage{xcolor}
\newtheorem{theorem}{Theorem}
\newcommand{\T}{\mathbb T}
\newcommand{\cD}{\mathcal D}
\newcommand{\cM}{\mathcal M}
\newcommand{\e}{\mathrm e}
\title{Improved Weyl bounds on short intervals}
\author{Xiyu Hu}
\date{Source: arXiv:2609.02478v1 (CC BY 4.0)}
\begin{document}
\maketitle

\noindent\emph{Source and licence.} Improved Weyl bounds on short intervals. Version arXiv:2609.02478v1 (CC BY 4.0), distributed by arXiv under the Creative Commons Attribution 4.0 International licence, http://creativecommons.org/licenses/by/4.0/, as recorded in the arXiv licence metadata for this exact version and shown on the abstract page https://arxiv.org/abs/2609.02478v1. Full licence notice: \texttt{licenses/arxiv-2609.02478-CC-BY-4.0.txt}.\par\smallskip

\noindent\textit{Editorial scope.} Counted unit: the article's theorem thm:maximal-fibre-moment (source0.tex lines 665--675). The article's complete original proof of it is delivered with it (source0.tex lines 677--720, one \texttt{proof} environment of 44 lines). No mathematics is written, repaired or re-derived: the statement and the proof are the article's own lines, in the article's own notation, and no retained line is substituted or re-flowed.  Retained uncounted as the source-local context the proof needs: lines 610--614.  Nothing was omitted from any retained span, so the package carries no omission record.  No retained line contains a citation, so the wrapper carries no bibliography and no bibliography entry.  No retained line contains a figure environment, an image-inclusion command or drawing code, so no image is delivered and the package declares no asset dependency.  Editorial formatting only, in addition: an AMS \texttt{amsart} preamble that restates the article's own declarations and macros used by the retained lines, namely the environments \texttt{theorem} on separate counters, together with the standard packages those lines need.  The retained environments print in this document as Theorem 1 in order of appearance, whereas the article prints the same items with the numbering of its own full document.  The complete native source of the article as distributed in the arXiv e-print for this exact version is bundled unchanged as \texttt{sources/209203/source0.tex}.
\begin{equation}\label{eq:critical-vmvt}
  s=\frac{r(r+1)}2,
  \qquad
  J_{s,r}(N)\ll_{r,\varepsilon}N^{s+\varepsilon}.
\end{equation}

\begin{theorem}\label{thm:maximal-fibre-moment}
For every fixed $k\ge3$, every $\varepsilon>0$, and every $\theta\in\T$,
\begin{equation}\label{eq:maximal-fibre-moment}
  \int_{\T^{k-1}}
  \cM_{\theta,N}(\boldsymbol\beta)^K
  d\boldsymbol\beta
  \ll_{k,\varepsilon}
  N^{K/2+\varepsilon}.
\end{equation}
The implied constant is uniform in $\theta$.
\end{theorem}

\begin{proof}
Let $s=K/2=k(k-1)/2$, which is the critical exponent for the Vinogradov system of degree $k-1$.  First consider an integer interval $J$ of length $R\le N$.  Expanding the $2s$th moment and integrating in $\beta_1,\ldots,\beta_{k-1}$ gives
\begin{align*}
  &\int_{\T^{k-1}}
  \left|\sum_{n\in J}
  \e\left(\theta n^k+\sum_{j=1}^{k-1}\beta_jn^j\right)
  \right|^{2s}
  d\boldsymbol\beta\\
  &\qquad=
  \sum_{\substack{x_1,\ldots,x_s,y_1,\ldots,y_s\in J\\
  \sum_i x_i^j=\sum_i y_i^j\ (1\le j\le k-1)}}
  \e\left(\theta\left(\sum_i x_i^k-\sum_i y_i^k\right)\right).
\end{align*}
Taking absolute values and translating $J$ to an interval beginning at one, the right-hand side is at most $J_{s,k-1}(R)$.  Translation preserves the system of equal power sums because the numbers of variables on the two sides are equal.  Hence \eqref{eq:critical-vmvt} gives
\begin{equation}\label{eq:interval-fibre-moment}
  \int_{\T^{k-1}}
  \left|\sum_{n\in J}
  \e\left(\theta n^k+\sum_{j=1}^{k-1}\beta_jn^j\right)
  \right|^{K}
  d\boldsymbol\beta
  \ll_{k,\varepsilon}R^{K/2+\varepsilon}.
\end{equation}

Every initial interval $[1,L]$ is the disjoint union of at most $1+\log_2N$ dyadic intervals from a fixed dyadic grid $\cD_N$.  Therefore
\[
  \cM_{\theta,N}(\boldsymbol\beta)^K
  \ll_k
  (\log N)^{K-1}
  \sum_{J\in\cD_N}
  \left|\sum_{n\in J}
  \e\left(\theta n^k+\sum_{j=1}^{k-1}\beta_jn^j\right)
  \right|^K.
\]
At dyadic scale $R$ there are $O(N/R)$ intervals.  Integrating and using \eqref{eq:interval-fibre-moment}, we obtain
\[
  \int_{\T^{k-1}}\cM_{\theta,N}^K
  \ll_{k,\varepsilon}
  N^\varepsilon
  \sum_{R\text{ dyadic}\le N}
  \frac NR R^{K/2+\varepsilon}
  \ll_{k,\varepsilon}N^{K/2+2\varepsilon}.
\]
Renaming $2\varepsilon$ as $\varepsilon$ proves the theorem.
\end{proof}

\end{document}
